\subsection{IDEA}
\subsubsection{IDEA简介}
IDEA（International Data Encryption Algorithm）是1991年由James Massey和Xuejia Lai提出的分组密码算法，
原名为IPES（International Pseudo-Efficient Scheme）。IDEA使用128位密钥，分组长度为64位，采用8轮加密。

\subsubsection{轮函数}
IDEA轮函数不是Feistel结构，其中主要有三种运算：异或、模$2^{16}$加法和模$2^{16}+1$乘法。
每轮将输入分为4个16位分组，与6个16位子密钥进行运算。
\begin{figure}[H]
	\centering
	\includegraphics[width=0.4\textwidth, keepaspectratio]{3-5.pdf}
	\caption{IDEA轮结构}
\end{figure}

\subsubsection{输出变换}
8轮加密后，IDEA使用4个子密钥进行输出变换，注意$M_2$与$M_3$交换了位置。
\begin{figure}[H]
	\centering
	\includegraphics[width=0.5\textwidth, keepaspectratio]{3-6.pdf}
	\caption{输出变换}
\end{figure}

\subsubsection{子密钥产生}
IDEA的子密钥由128位密钥生成，分为52个16位子密钥。子密钥生成过程如下：
\begin{enumerate}
    \item 将128位密钥分为8个16位分组。
    \item 对每个分组进行循环左移25位，生成子密钥。
    \item 生成的子密钥按顺序使用。
\end{enumerate} 